\chapter{Cohomology, Levi Decomposition, and Ado}
\label{ch:cohomology-levi-ado}

An extension remembers an action as well as a kernel. Throughout the first
section, \(\mathfrak g\) is a Lie algebra over a field \(k\), and \(M\) is
a specified left \(\mathfrak g\)-module. The Chevalley--Eilenberg convention is
\[
 (db)(x,y)=x b(y)-y b(x)-b([x,y]).
\]
We use the exterior-power complex, its coefficient long exact sequence,
and its comparison with \(\operatorname{Ext}_{U(\mathfrak g)}\) proved in
\cite[Chapter 6, A06--A08]{homological}. In particular, all coefficient
modules in an application of Whitehead's lemmas below are finite dimensional.

\section{Extensions and deformations}

\begin{theorem}[Abelian extensions with a specified action]\label{res:I01}
Equivalence classes of extensions
\[
 0\longrightarrow M\longrightarrow E\longrightarrow\mathfrak g
 \longrightarrow0
\]
whose kernel is abelian and whose induced action on it is the given one
are parametrized by \(H^2(\mathfrak g,M)\). Equivalence induces the identity
on both \(M\) and \(\mathfrak g\). Central extensions are the special case
of the trivial action. An extension has a Lie algebra section if and only
if its class is zero.
\end{theorem}
\begin{proof}
Choose a linear section \(s\). The error
\(\omega(x,y)=[s(x),s(y)]-s([x,y])\) lies in \(M\).
Under the resulting identification \(E=\mathfrak g\oplus M\), the bracket is
\[
 [(x,m),(y,n)]=([x,y],\,xn-ym+\omega(x,y)).
 \tag{9.1}\label{eq:abelian-extension}
\]
Jacobi on three elements of \(\mathfrak g\) has \(M\)-component
\[
 x\omega(y,z)-y\omega(x,z)+z\omega(x,y)
 -\omega([x,y],z)+\omega([x,z],y)-\omega([y,z],x).
\]
Thus it is precisely \(d\omega=0\). Jacobi with one entry in \(M\)
is the module identity, and with two entries in \(M\) it is zero.
Consequently every \(2\)-cocycle produces an extension by
\eqref{eq:abelian-extension}.

Changing the section to \(s+b\) changes \(\omega\) to \(\omega+db\).
Conversely any linear equivalence of extensions has the form
\((x,m)\mapsto(x,m+b(x))\). Substitution in \eqref{eq:abelian-extension}
shows that it preserves brackets exactly when the two cocycles differ
by the corresponding coboundary. This proves both directions of the
classification. The section \(x\mapsto(x,0)\) is a Lie section exactly
when \(\omega=0\); changing sections proves the last assertion.
See \cite[Section 7.6]{weibel} and \cite[Proposition 48.3]{etingof}.
\end{proof}

\begin{proposition}[First-order bracket deformations]\label{res:I02}
Fix a bracket \(\mu(x,y)=[x,y]\). Alternating brackets
\(\mu+t c\) on \(\mathfrak g\otimes_k k[t]/(t^2)\), up to changes of
coordinates \(1+t a\), are classified by
\(H^2(\mathfrak g,\mathfrak g)\), for the adjoint action.
\end{proposition}
\begin{proof}
The coefficient of \(t\) in Jacobi is
\[
 \sum_{\mathrm{cyc}}\bigl([c(x,y),z]+c([x,y],z)\bigr)=-dc(x,y,z).
\]
It vanishes exactly when \(c\) is a cocycle. Pulling back the bracket
by \(T=1+t a\), with inverse \(1-t a\), replaces \(c\) by
\[
 c(x,y)+[a(x),y]+[x,a(y)]-a([x,y])=c(x,y)+da(x,y).
\]
Every allowed change of coordinates has this form. Hence its orbits
are exactly cocycles modulo coboundaries. The action of the kernel
\(t\mathfrak g\) is adjoint; it is not the trivial action of a central
extension. This also recovers the assertion from
Theorem~\ref{res:I01}. Compare \cite[Proposition 48.5]{etingof}.
\end{proof}

\begin{proposition}[Obstructions and formal rigidity]\label{res:I03}
Suppose \(k\) has characteristic zero. A bracket
\(\mu_t=\sum_{i=0}^{N-1}t^i c_i\), where \(c_0=\mu\), that satisfies
Jacobi modulo \(t^N\) extends modulo \(t^{N+1}\) exactly when the class
of
\[
 O_N(x,y,z)=
 \sum_{\substack{i+j=N\\i,j>0}}
 \sum_{\mathrm{cyc}}c_i(c_j(x,y),z)
 \tag{9.2}\label{eq:lie-obstruction}
\]
vanishes in \(H^3(\mathfrak g,\mathfrak g)\). In particular
\(H^2(\mathfrak g,\mathfrak g)=0\) implies that every one-parameter
formal deformation is formally equivalent to the constant bracket.
\end{proposition}
\begin{proof}
The terms involving the unknown \(c_N\) are \(-dc_N\), by the
calculation in Proposition~\ref{res:I02}. The equation to solve is
therefore \(dc_N=O_N\).
Here is the closure check on the obstruction. On alternating maps use
the insertion operation
\[
 (f\circ h)(x_1,\ldots,x_{p+q+1})
 =\sum_{\sigma\in\operatorname{Sh}(q+1,p)}
 \operatorname{sgn}(\sigma)
 f\bigl(h(x_{\sigma(1)},\ldots,x_{\sigma(q+1)}),
          x_{\sigma(q+2)},\ldots,x_{\sigma(p+q+1)}\bigr),
\]
where \(f\) has arity \(p+1\) and \(h\) has arity \(q+1\).
Its commutator
\([f,h]_{\mathrm{NR}}=f\circ h-(-1)^{pq}h\circ f\)
satisfies graded Jacobi. In fact, for \(\ell\) of arity
\(r+1\), the shuffle expansion gives the associator identity
\[
 (f\circ h)\circ \ell-f\circ(h\circ \ell)
   =(-1)^{qr}\bigl((f\circ \ell)\circ h-f\circ(\ell\circ h)\bigr).
\]
On each side the nested-block terms cancel. The remaining terms
correspond to disjoint input sets for the two insertions.
Reordering those blocks swaps \(q+1\) and \(r+1\) inputs,
and their outputs occupy two slots of the alternating map \(f\).
Moving each output past the other block before the final slot
exchange gives the additional signs
\((-1)^{q+1}\), \((-1)^{r+1}\) and \(-1\); their product
with \((-1)^{(q+1)(r+1)}\) is \((-1)^{qr}\).
Expanding the three nested commutators in graded Jacobi pairs
these associators with precisely this sign, so they cancel.
For a binary alternating \(\nu\), its Jacobiator is
\(\tfrac12[\nu,\nu]_{\mathrm{NR}}=\nu\circ\nu\).
The same associator identity with \(f=h=\ell=\nu\) gives
\([\nu,[\nu,\nu]_{\mathrm{NR}}]_{\mathrm{NR}}=0\).
Apply this to the truncated polynomial \(\mu_t\). Its Jacobiator has
no coefficients below \(N\); the coefficient of \(t^N\) in this
identity is \([\mu,O_N]_{\mathrm{NR}}=0\), which is \(dO_N=0\),
up to the fixed overall sign in the exterior-complex convention.
Thus the displayed equation is solvable exactly when its cohomology
class vanishes. Altering a proposed \(c_N\) by a coboundary changes
coordinates and not the obstruction class.

For the rigidity assertion, suppose the first nonzero coefficient is
\(t^m c_m\). It is a \(2\)-cocycle. Write \(c_m=da_m\), and pull back
by \(1-t^m a_m\) to kill it without changing lower coefficients.
Repeat at degrees \(m+1,m+2,\ldots\). The product of these changes
converges in the \(t\)-adic topology: each coefficient is altered
only finitely often, and the constant term remains the identity,
so its inverse exists by formal recursion. The limiting change
removes every positive coefficient.
This proves formal, not analytic-convergence, rigidity.
See \cite[Section 48.1]{etingof} and \cite[Sections 1 and 3]{fialowski}.
\end{proof}

\section{Whitehead's second lemma and Levi factors}

\begin{theorem}[Second Whitehead lemma for trivial coefficients]\label{res:I04}
If \(\mathfrak g\) is finite dimensional and semisimple over a
characteristic-zero field, then \(H^2(\mathfrak g,k)=0\).
\end{theorem}
\begin{proof}
The Killing form \(B\) is nondegenerate by Theorem~\ref{res:B05}.
Given an alternating \(2\)-cocycle \(\omega\), define \(D\) uniquely by
\(\omega(x,y)=B(Dx,y)\). Alternation says that \(D\) is skew-adjoint
for \(B\). Pairing
\(D[x,y]-[Dx,y]-[x,Dy]\) with \(z\), using invariance of \(B\),
and using skew-adjointness, gives
\[
 \omega([x,y],z)+\omega([y,z],x)+\omega([z,x],y)=0.
\]
Thus \(D\) is a derivation. By Theorem~\ref{res:B07}, \(D=\operatorname{ad}a\).
It follows that
\(\omega(x,y)=B(a,[x,y])=d\ell(x,y)\) for
\(\ell(x)=-B(a,x)\). Every cocycle is a coboundary.
This is the trivial-coefficient case of
\cite[Corollary 7.8.12]{weibel}.
\end{proof}

\begin{theorem}[Second Whitehead lemma with coefficients]\label{res:I05}
For a finite-dimensional semisimple characteristic-zero Lie algebra
\(\mathfrak g\) and a finite-dimensional module \(M\),
\[
 H^2(\mathfrak g,M)=0.
\]
Every abelian extension of \(\mathfrak g\) by such a specified module
has a Lie algebra section.
\end{theorem}
\begin{proof}
By Weyl's theorem, Theorem~\ref{res:C05}, \(M\) is a finite direct
sum of simple modules. The exterior-complex differential respects
this finite direct sum. A nontrivial simple summand has zero
cohomology in every degree by the Casimir homotopy in
Theorem~\ref{res:C03}; a trivial simple summand is covered by
Theorem~\ref{res:I04}. Thus degree two vanishes for their direct sum.
The extension assertion follows from Theorem~\ref{res:I01}.
This argument retains finite-dimensionality and does not claim
that all higher cohomology with trivial coefficients vanishes.
See \cite[Corollary 7.8.12]{weibel}.
\end{proof}

\begin{theorem}[Levi decomposition]\label{res:I06}
Let \(\mathfrak g\) be a finite-dimensional Lie algebra over a
characteristic-zero field, with solvable radical \(\mathfrak r\).
There is a semisimple subalgebra \(\mathfrak s\) such that
\[
 \mathfrak g=\mathfrak s\ltimes\mathfrak r,\qquad
 \mathfrak s\xrightarrow{\;\sim\;}\mathfrak g/\mathfrak r.
\]
\end{theorem}
\begin{proof}
Induct on \(\dim\mathfrak r\). If it is zero, take \(\mathfrak s=\mathfrak g\).
Otherwise take the last nonzero term \(\mathfrak a\) of the derived
series of \(\mathfrak r\); if \(\mathfrak r\) is abelian, take
\(\mathfrak a=\mathfrak r\). This is a nonzero abelian characteristic
ideal of \(\mathfrak r\), hence an ideal of \(\mathfrak g\).
The radical of \(\mathfrak g/\mathfrak a\) is
\(\mathfrak r/\mathfrak a\): the latter is solvable, and the quotient
by it is semisimple. Induction gives a Levi subalgebra
\(\overline{\mathfrak s}\) in this quotient.

Let \(E\) be its inverse image in \(\mathfrak g\).
Then \(0\to\mathfrak a\to E\to\overline{\mathfrak s}\to0\)
is an abelian extension, with its actual adjoint action on
\(\mathfrak a\). By Theorem~\ref{res:I05} its class is zero.
A Lie section has image \(\mathfrak s\).
Its intersection with \(\mathfrak r\) is zero because
\(\overline{\mathfrak s}\cap(\mathfrak r/\mathfrak a)=0\) and the
section has no intersection with \(\mathfrak a\).
Moreover its map to \(\mathfrak g/\mathfrak r\) is surjective,
so dimensions give the required vector-space direct sum.
The radical is an ideal, so this is the stated semidirect product.
See \cite[Theorem 7.8.13]{weibel} and
\cite[Theorem 48.9]{etingof}.
\end{proof}

\begin{theorem}[Conjugacy of Levi factors; reductive algebras]\label{res:I07}
Any two Levi factors are conjugate by
\(\exp(\operatorname{ad}z)\) for an element \(z\) of the nilradical
\(\mathfrak n\). The exponential is a finite sum.
A finite-dimensional characteristic-zero Lie algebra whose radical
is its center is the direct sum of a semisimple algebra and its center.
\end{theorem}
\begin{proof}
Fix one Levi factor \(\mathfrak s\) and write the second as
\(\{x+f(x):x\in\mathfrak s\}\), where \(f:\mathfrak s\to\mathfrak r\)
is linear. By Theorem~\ref{res:B08},
\([\mathfrak g,\mathfrak r]\subseteq\mathfrak n\).
Modulo \(\mathfrak n\), therefore,
\(\mathfrak g/\mathfrak n=\mathfrak s\oplus(\mathfrak r/\mathfrak n)\)
has central second summand. Since \(\mathfrak s=[\mathfrak s,\mathfrak s]\),
the homomorphism into that central summand is zero. Thus \(f\)
already takes values in \(\mathfrak n\).

Write \(\gamma_1\mathfrak n=\mathfrak n\) and
\(\gamma_{j+1}\mathfrak n=[\mathfrak n,\gamma_j\mathfrak n]\).
Suppose a conjugate second factor has discrepancy in
\(\gamma_j\mathfrak n\). In
\(\gamma_j\mathfrak n/\gamma_{j+1}\mathfrak n\) the equation for its
bracket is
\[
 f([x,y])=x f(y)-y f(x);
\]
the additional term \([f(x),f(y)]\) belongs to
\(\gamma_{2j}\mathfrak n\subseteq\gamma_{j+1}\mathfrak n\).
By Theorem~\ref{res:C04} there is \(v_j\in\gamma_j\mathfrak n\)
such that \(f(x)=[x,v_j]\) modulo \(\gamma_{j+1}\mathfrak n\).
Conjugating the second factor by \(\exp(\operatorname{ad}v_j)\)
changes the discrepancy by \([v_j,x]\), and kills it modulo
\(\gamma_{j+1}\mathfrak n\). Higher exponential terms have values
in that next term as well.

The lower central series terminates. Also
\(\operatorname{ad}v_j\) is nilpotent on all of \(\mathfrak g\):
its first application lands in the nilpotent ideal \(\mathfrak n\),
and its subsequent applications follow that ideal's lower central
series. Thus all the exponentials used are finite. Their product
is \(\exp(\operatorname{ad}z)\) for \(z\in\mathfrak n\), by the
finite BCH identity of Theorem~\ref{res:A09}. It sends the
second factor to the first.

For the last assertion, Levi decomposition writes the algebra as
\(\mathfrak s\ltimes Z(\mathfrak g)\). The action on the center is
zero, hence the sum is direct. Conversely such a direct sum has
exactly its abelian central summand as radical. This is the
meaning of \emph{reductive} used here; it is not an assertion that
every representation of a reductive algebra is completely reducible.
The conjugacy statement expands the Malcev observation following
\cite[Theorem 7.8.13]{weibel}.
\end{proof}

\section{Finite faithful representations}

\begin{proposition}[A derivation-stable nilpotent truncation]\label{res:I08}
Every finite-dimensional nilpotent Lie algebra \(\mathfrak n\)
has a finite-dimensional faithful representation by nilpotent
operators. The construction can be made compatible with every
derivation of \(\mathfrak n\).
\end{proposition}
\begin{proof}
Choose a basis adapted to the lower central series. Give a basis
element weight \(j\) if it belongs to
\(\gamma_j\mathfrak n\) but not to \(\gamma_{j+1}\mathfrak n\).
For an ordered PBW monomial give its weight as the sum of the
weights of its factors. Let \(F_{\geq N}U(\mathfrak n)\) be the span
of monomials of weight at least \(N\).

This is a two-sided ideal. In fact, inserting a new factor and
reordering by PBW either preserves the sum of weights or replaces
two factors by a bracket. The latter has weight at least their
sum because
\([\gamma_i\mathfrak n,\gamma_j\mathfrak n]\subseteq
\gamma_{i+j}\mathfrak n\). No resulting term has smaller weight.
Every derivation \(\delta\) preserves each \(\gamma_j\mathfrak n\),
as follows by induction from the derivation identity. Its extension
to \(U(\mathfrak n)\), obtained by applying it in every factor, also
never decreases weight. Thus it preserves this ideal.

The quotient \(V_N=U(\mathfrak n)/F_{\geq N}\) is finite dimensional:
the positive integer weights bound each exponent in a monomial of
weight less than \(N\). If \(\gamma_{c+1}\mathfrak n=0\), take
\(N>c\). The distinct PBW vectors \(1\) and all the individual
basis vectors survive and remain independent. Hence left
multiplication \(L_x\) is faithful on \(\mathfrak n\), since
\(L_x(1)=x\). It raises weight by at least one, so \(L_x^N=0\).
Finally every \(\delta\) acts on \(V_N\), kills \(1\), and satisfies
\([D_\delta,L_x]=L_{\delta x}\).
This is a PBW proof of the faithful nilpotent-module conclusion
of \cite[Proposition 50.5 and Corollary 50.7]{etingof}.
\end{proof}

\begin{proposition}[A reductive--nilpotent enlargement]\label{res:I09}
Every finite-dimensional complex Lie algebra embeds into an algebra
\[
 \widehat{\mathfrak g}=\mathfrak k\ltimes\mathfrak n,\qquad
 \mathfrak k=\mathfrak s\oplus\mathfrak t,
\]
where \(\mathfrak s\) is semisimple, \(\mathfrak t\) is abelian
and acts semisimply, and \(\mathfrak n\) is nilpotent.
The central toral actions introduced in the construction have
integer weights in a chosen basis of \(\mathfrak t\).
\end{proposition}
\begin{proof}
We give an entirely Lie-algebraic version of the reduction in
\cite[Proposition 50.3]{etingof}. Inductively maintain a decomposition
\(\mathfrak k\ltimes\mathfrak a\), with \(\mathfrak a\) solvable,
\(\mathfrak k=\mathfrak s\oplus\mathfrak t\) reductive, and
\(\mathfrak t\) acting diagonally on \(\mathfrak a\).
Such an algebra starts the induction: take a Levi factor for
\(\mathfrak k=\mathfrak s\), and \(\mathfrak a=\mathfrak r\).
Weyl's theorem makes the initial \(\mathfrak k\)-module
\(\mathfrak a\) completely reducible. More generally, a module
on which the center acts diagonally is a direct sum of its central
weight spaces; on each, Weyl's theorem for \(\mathfrak s\)
gives complete reducibility.

Let \(\mathfrak n_0\) be the nilradical of \(\mathfrak a\).
Use the nonnegative integer \(q=\dim(\mathfrak a/\mathfrak n_0)\)
as decreasing measure. If \(q=0\), we are done.
Every derivation of \(\mathfrak a\) acts trivially on
\(\mathfrak a/\mathfrak n_0\), by Theorem~\ref{res:B08}.
Complete reducibility therefore supplies a
\(\mathfrak k\)-invariant element \(d\in\mathfrak a\)
whose image in that quotient is nonzero. Choose a linear
functional \(\chi\) on \(\mathfrak a/\mathfrak n_0\) with
\(\chi(d)=1\), and regard it as a functional on \(\mathfrak a\).
It is a character because
\([\mathfrak a,\mathfrak a]\subseteq\mathfrak n_0\).

Decompose \(\mathfrak a=\bigoplus_\beta\mathfrak a_{[\beta]}\)
into generalized eigenspaces of \(D=\operatorname{ad}d\).
The derivation identity gives
\[
 [\mathfrak a_{[\beta]},\mathfrak a_{[\gamma]}]
 \subseteq\mathfrak a_{[\beta+\gamma]}.
 \tag{9.3}\label{eq:additive-generalized-weights}
\]
Indeed \((D-(\beta+\gamma))^m[x,y]\) is the binomial sum of brackets
of \((D-\beta)^i x\) and \((D-\gamma)^{m-i}y\), and is zero for
large \(m\). The finitely many eigenvalues generate a free
abelian subgroup \(\Gamma\subset\mathbb C\); choose a
\(\mathbb Z\)-basis \(\alpha_1,\ldots,\alpha_m\).
Write \(\beta=\sum_j b_{\beta j}\alpha_j\), with integers
\(b_{\beta j}\).
An abelian algebra \(\mathfrak t_1=\mathbb C^m\) acts on
\(\mathfrak a_{[\beta]}\) by
\((z_j)\mapsto\sum_j b_{\beta j}z_j\).
Equation~\eqref{eq:additive-generalized-weights} proves that these
operators are derivations. They commute with \(\mathfrak k\),
because \(d\) is \(\mathfrak k\)-invariant.

For \(\alpha=(\alpha_1,\ldots,\alpha_m)\in\mathfrak t_1\),
the action of \(\alpha\) is the semisimple part \(S\) of \(D\).
Both \(\mathfrak n_0\) and \(\ker\chi\) are stable under
\(\mathfrak t_1\). For the first, \(D\)-stability gives a
decomposition into generalized eigenspaces; for the second,
\(D\) acts trivially on \(\mathfrak a/\mathfrak n_0\), so that
only weight zero occurs in this quotient.
Also \(\mathfrak t_1\) kills \(d\), since \(D d=0\).

In \((\mathfrak k\oplus\mathfrak t_1)\ltimes\mathfrak a\), put
\[
 y=d-\alpha,\qquad
 \mathfrak a'=\ker\chi\oplus\mathbb C y.
\]
It is an ideal complementary to
\(\mathfrak k'=\mathfrak k\oplus\mathfrak t_1\).
For example \([y,x]=(D-S)x\) for \(x\in\ker\chi\), and this
lies in \(\mathfrak n_0\), since both \(D\) and \(S\) act
trivially on \(\mathfrak a/\mathfrak n_0\).
All other required brackets follow from the stability just proved.
The new algebra contains the old \(\mathfrak k\ltimes\mathfrak a\);
we have only adjoined \(\mathfrak t_1\).

The ideal \(\mathfrak m=\mathfrak n_0+\mathbb C y\) of
\(\mathfrak a'\) is nilpotent. To check this carefully, every
\(x\in\mathfrak n_0\) acts nilpotently on \(\mathfrak a'\):
its first bracket lands in \(\mathfrak n_0\), where repeated
bracketing terminates. The element \(y\) acts on \(\ker\chi\)
by the nilpotent operator \(D-S\) and kills itself.
Lie's theorem simultaneously upper triangularizes the adjoint
representation of the solvable algebra \(\mathfrak m\) on
\(\mathfrak a'\). Its diagonal characters vanish on
\(\mathfrak n_0\) and on \(y\), by the two nilpotency observations,
so every element has zero diagonal. Engel's theorem gives the
claim. Hence the new nilradical contains \(\mathfrak m\), and
\[
 \dim(\mathfrak a'/\operatorname{nilrad}\mathfrak a')
 \leq \dim(\mathfrak a/\mathfrak n_0)-1.
\]

Finally the central torus of \(\mathfrak k'\) acts diagonally
on \(\mathfrak a'\): it does so on the old generalized
eigenspace blocks, and it kills the extra line \(\mathbb C y\).
Its new weights are the integers \(b_{\beta j}\).
Thus complete reducibility needed in the next step is still
available. The integer \(q\) decreases until it is zero.
This proves the enlargement without assuming that arbitrary
modules of a reductive algebra are semisimple, or assuming an
integration theorem for an algebraic group.
\end{proof}

\begin{proposition}[The semidirect-product representation]\label{res:I10}
If \(\mathfrak k\) is a finite-dimensional complex reductive algebra
and it acts by derivations on a finite-dimensional nilpotent
algebra \(\mathfrak n\), then \(\mathfrak k\ltimes\mathfrak n\)
has a finite-dimensional faithful module.
\end{proposition}
\begin{proof}
Take \(V_N\) from Proposition~\ref{res:I08}. The formula
\[
 (a,x)\longmapsto D_a+L_x
\]
is a representation: \([D_a,D_b]=D_{[a,b]}\),
\([D_a,L_y]=L_{a y}\), and \([L_x,L_y]=L_{[x,y]}\)
give exactly the semidirect-product bracket.
If \(D_a+L_x\) is zero, applying it to \(1\) gives \(x=0\).
Its restriction to the surviving copy of \(\mathfrak n\)
then says that \(a\) acts as the zero derivation on \(\mathfrak n\).

To detect this possible remaining kernel, take a faithful module
\(W\) of \(\mathfrak k=\mathfrak s\oplus\mathfrak t\).
The adjoint module of \(\mathfrak s\) is faithful because its center
is zero. If \(t_1,\ldots,t_m\) is a basis of \(\mathfrak t\),
let \(t_j\) act on \(\mathbb C^m\) by the diagonal matrix with
its only nonzero entry \(1\) in position \(j\); let \(\mathfrak s\)
act trivially on that summand and \(\mathfrak t\) trivially on
the adjoint summand. Their direct sum is the required \(W\).
Inflate it along the projection
\(\mathfrak k\ltimes\mathfrak n\to\mathfrak k\).
On \(V_N\oplus W\), the kernel condition first gives \(x=0\)
and then \(a=0\). This also covers \(\mathfrak n=0\) and pure
reductive direct factors; no faithfulness of the action of
\(\mathfrak k\) on \(\mathfrak n\) has been presumed.
Compare the final construction in \cite[Theorem 50.8]{etingof}.
\end{proof}

\begin{theorem}[Ado over the complex and real fields]\label{res:I11}
Every finite-dimensional complex or real Lie algebra admits an
injective homomorphism into \(\mathfrak{gl}(V)\) for a
finite-dimensional vector space over the same field.
\end{theorem}
\begin{proof}
Over \(\mathbb C\), embed into the enlargement of
Proposition~\ref{res:I09}, apply Proposition~\ref{res:I10},
and restrict the faithful representation to the original subalgebra.
For a real algebra, apply the complex assertion to
\(\mathfrak g\otimes_{\mathbb R}\mathbb C\) and regard its
representation as a real vector space. The original map
\(\mathfrak g\to\mathfrak g\otimes_{\mathbb R}\mathbb C\)
is injective, so its action remains faithful. The zero algebra
may act on the zero-dimensional space, or trivially on a line.
The arbitrary characteristic-zero version is proved in
Proposition~\ref{res:J02}, with its descent step included there.
\end{proof}
